\section{Introduction}

This is a small, hand-written Carve document. It exercises the
constructs that survive a round-trip through Pandoc into \emph{every}
writer: emphasis, \textbf{strong} text, links, lists, tables, footnotes,
math, and raw passthrough.

Carve swaps the Markdown emphasis delimiters: \texttt{/italic/} is
emphasis and \texttt{*bold*} is strong. A footnote reference looks like
this.\footnote{Footnotes become Pandoc \texttt{Note} nodes and render
  natively in LaTeX, Typst, DOCX, and the rest.}

\subsection{Lists}

\begin{itemize}
\tightlist
\item
  First item
\item
  Second item, with \emph{emphasis}

  \begin{itemize}
  \tightlist
  \item
    A nested item
  \end{itemize}
\item
  Third item
\end{itemize}

\begin{enumerate}
\def\labelenumi{\arabic{enumi}.}
\tightlist
\item
  Ordered one
\item
  Ordered two
\end{enumerate}

\subsection{A table}

\begin{longtable}[]{@{}lll@{}}
\toprule\noalign{}
Format & Binary? & Text golden? \\
\midrule\noalign{}
\endhead
\bottomrule\noalign{}
\endlastfoot
LaTeX & no & yes \\
Typst & no & yes \\
DOCX & yes & no \\
\end{longtable}

\subsection{Math and raw passthrough}

Inline math: \(e^{i\pi} + 1 = 0\).

A target-routed raw span reaches only the LaTeX writer and is dropped by
every other format: \LaTeX. This is the pandoc-Markdown raw-span
semantics that Carve's \texttt{\{=format\}} syntax was born from.

\begin{quote}
A blockquote closes the tour. Every node above maps to a Pandoc AST
node, so \texttt{pandoc} can then emit LaTeX, Typst, RST, DOCX, and
dozens more.
\end{quote}
